\author{OTS}
\documentclass[fontsize=14pt]{../../inc/mybib}

\setincpath{../../inc/}

\graphicspath{{../../assets/images/}}

\newcommand{\Name}{Fredy}
\setheaderlogo{mnr-green.png}

% ensure scrlayer-scrpage has sufficient footheight
\setlength{\footheight}{30.4pt}

\begin{document}
\begin{bibelbox}{SCHL}{Kol}{3:15-17}
    Lasst das Wort des Christus reichlich in euch wohnen in aller Weisheit; lehrt und ermahnt einander und singt mit Psalmen und Lobgesängen und geistlichen Liedern dem Herrn lieblich in eurem Herzen. Und was immer ihr tut in Wort und Werk, das tut alles im Namen des Herrn Jesus und dankt Gott, dem Vater, durch ihn.
\end{bibelbox}
\section{Begrüssung}
Ich möchte auch euch alle recht herzlich zum Gottesdienst begrüssen. Schön, dass ihr gekommen seid, um unserem \herr N zu danken, ihn zu loben und zu preisen. Nick und Erika lass sich entschuldigen, dass sie heute nicht hier sein können. Nick musste kurzfristig nach Ungarn fahren, weil eine Freundin von ihm verstorben ist. 

Auch dich \Name{} will ich herzlich begrüssen. Du bist heute das erste Mal hier bei uns und wir freuen uns sehr dich hier begrüssen zu können.
\begin{itemize}
    \item[] \beten[b, p]{Beten zur Begrüssung}
    \item[] und anschliessend singen wir zusammen das Lied
    \item[] \lied[b]{Du bist das Licht der Welt.}
\end{itemize}

\section{Informationen}
\begin{itemize}
    \item \info[b, p]{Bibel und Gebetsabend:} Do., 04. Juni 2026 20:00 Uhr Bibel und Gebetsabend mit Markus 8,10-13 Samuel Rindlisbacher.\\
    Vom 22.6. bis 19.7.2026 sind Janina und ich auf der Alp Erner Galen zum Schafe hüten. Die Gottesdienste während der Zeit werden normal weiter geführt, aber am Donnerstagabend, findet nur ein Gebetsabend statt ohne Liveübertragung. Wenn jemand kommt, der die Liveübertragung auf den Beamer projizieren kann, dann findet die Übertragung statt. Wenn niemand da ist, der weiss wie das funktioniert, findet nur das Gebet statt.
    \item \info[b, p]{Nächster Gottesdienst:} So, 14. Juni 2026 Lothar und Israel Reise
    \item \info[b, p]{Kollekte:} Die Kollekte wird hinten in der grünen Box gesammelt und wird vollumfänglich für den Bau dieser Gemeinde eingesetzt.
\end{itemize}

\section{Input}
\begin{spacing}{1.5}
   \begin{block}[Gottesnamen]
    Machen wir weiter mit den Namen Gottes. Wir haben die drei Hauptnamen YHWH (Jahwe oder Jehova), Adonai und Elohim. In der Bibel finden wir aus diesen Hauptnamen diverse zusammengesetzte Namen, die oft die Eigenschaften unseres grossen Gottes beschreiben. Zwei haben wir uns schon angeschaut, Jahwe Zebaoth und Jahwe Adonai. Diese Namen haben wir angeschaut, weil die so häufig in der Bibel auftauche. Den Namen heute wollen wir anschauen, nicht weil er so häufig auftaucht, sondern weil er so häufig missbraucht wird. 

    Heute wollen wir uns Jahwe Rapha \enquote{der \herr{} dein Arzt} anschauen.  
    \end{block}
    \begin{block}[Der Herr ist dein Arzt]
        Dieser Name kommt nicht oft in der Bibel vor, aber er wird oft gebraucht und auch missbraucht. Als Eigennamen kommt er so nur in \bibleverse{IIMos}(15:26) vor:
        \begin{bibelbox}{SCHL}{IIMos}{15:26}
            und er sprach: Wenn du der Stimme des \herr N, deines Gottes, eifrig gehorchen wirst und tust, was ihm recht ist, und seine Gebote zu Ohren fasst und alle seine Satzungen hältst, so will ich keine der Krankheiten auf dich legen, die ich auf Ägypten gelegt habe; denn ich bin der \herr{} dein Arzt.
        \end{bibelbox}
        Lesen wir noch die Verse \bibleverse{Ps}(103:1-3)
        Dieser Vers wird so oft zitiert um zu sagen, dass man nur die Gebote von Gott halten soll und dann wird Gott einem alle Krankheiten geheilt. Oder anders gesagt, \enquote{wenn du nur genügend glaubst, wirst du nicht krank}, also \enquote{wenn du krank wirst, glaubst du zu wenig}. Es wird nicht so direkt gesagt, aber wenn du diesen Bibelvers einer kranken Person vorliest, dann wird dieser von dieser Person oft auch so verstanden. Fragt mal Fredy oder Nick, was es mit einem selber macht, wenn man diesen Vers hingeworfen bekommt.

        Ich bin $100\%$ überzeugt, dass Gott heilt, dass Gott geheilt hat und das Gott immer noch heilt. Wieso auch nicht? Er hat ja uns geschaffen, so kann er uns auch reparieren. 
        
        Man liest immer wieder Geschichten, dass Personen geheilt wurden. Aber ich habe noch nie gelesen, oder gehört, dass Gott nur geheilt hat, weil die Person alle Gebote gehalten hätte. Auch im alten Testament nicht. Zum Beispiel Naeman in \bibleverse{IKoen}(5:), er wurde von Gott geheilt, obwohl er nicht mal Jude war. Und im Neuen Testament wird dann berichtet, dass dieser der einzige Mensch in Israel war, der je vom Aussatz geheilt wurde. Jesus sagt in:
        \begin{bibelbox}{SCHL}{Lk}{4:27}
            und viele Aussätzige waren in Israel zur Zeit des Propheten Elisa; aber keiner von ihnen wurde gereinigt, sondern nur Naeman, der Syrer.
        \end{bibelbox}

        Was auch noch bemerkenswert ist, wenn Gott heilt, die Menschen wurden immer vollständig wieder hergestellt. Egal ob 40 Jahre gelähmt oder Tod, sie standen auf und konnten direkt wieder umherspazieren.

        Warum heilt Gott Menschen? Gott heilt Menschen zu seiner Ehre, es geschieht, dass Gott verherrlicht wird. Es geht nicht darum, in Talkshows aufzutreten oder mit Bücher Geld zu verdienen. Viele Geschichten, die im Umlauf sind, sollte man immer mit einer gewissen Skepsis anschauen. Auch bei sogenannten Nahtoderfahrungen. Am besten ist, man freut sich für diese Personen, die wissen am besten was und wie ihnen passiert ist. Verherrlichen Sie damit Jesus nicht, dann nehmt Abstand davon.

        Wir brauchen auch keine Heilungsveranstaltungen, wo Wunderheiler einem die Hand auflegen. Solche Heilungen sind nie von Gott. Es gibt auch keine medizinisch bezeugte Heilung aus solchen Veranstaltungen, aber was es Massenhaft gibt, sind Heiler die als Gauner überführt wurden.

        Zurück zum Namen. Der Name verspricht, dass Gott uns heilen wird, wenn wir seine Gebote halten. Im Alten Testament ist das irdische Schicksal des Volkes Israel am Halten der Gebote gehangen. Der sogenannte Mosebund, wir haben über diesen Bund im Winter gesprochen, ist der einzige Bund von Gott mit uns Menschen, der auf Bedingungen unserer seit's beruht. Jetzt sind wir aber im neuen Bund, in dem uns der Herr versprochen hat ein neues Herz zu geben. 

        Jesus kam auf die Welt und heilte viele Menschen. Vor allem um zu Beweisen, dass er der Messias ist. So lässt er Johannes dem Täufer ausrichten,
        \begin{bibelbox}{SCHL}{Matt}{11:5}
        Blinde werden sehend und Lahme gehend Aussätzige werden rein und Taube hören, Tote werden auferweckt, und Armen wird das Evangelium verkündigt.
        \end{bibelbox}

        War das, das Ziel Jesus, Menschen von ihren Krankheiten zu heilen? Bedeutet das \enquote{der Herr dein Arzt}. Nein! Der Herr dein Arzt heilt dich, in dem er sich selber schinden liess und für uns am Kreuz gestorben ist. Wenn es im alten Testament heisst, ich werde euch segnen, wenn ihr euch an meine Gebote hält, galt das nur für irdische Segnungen. Aber im Alten wie im Neuen Testament, werden wir durch den Glauben für ewig geheilt.

        Um geheilt zu werden, müssen wir auf das Kreuz blicken, auf den der für unsere Sünden gestorben ist. Er ist unser Arzt, der uns heilt. Nein, nicht von irdischen Krankheiten, sondern von der Krankheit der Sünde. Diese Krankheit ist wie Aussatz in der Ewigkeit und so aussätzig können wir nicht vor Gott treten. Für alle die hier auf Erden sich von ihren liebsten verabschieden mussten, wenn diese Gott als ihr Arzt von Herzen angenommen haben, sind sie für die Ewigkeit von diesem Aussatz geheilt und dürfen so vor das Angesicht Gottes durch Jesus Christus treten. Das ist der unser grösster Trost. 
        
        Der grösste Fitness Guru oder gesündeste Mensch der diesen Arzt nicht annimmt, wird auf ewig in der Ewigkeit mit dem Aussatz der Sünde behaftet sein und so für ewig von Gott getrennt.

        Lasst Gott euer Arzt sein. Dieser Arzt ist für unsere Heilung am Kreuz gestorben. Tragt das in die Welt hinaus, denn wir haben keinen Fachkräftemangel, weil Jesus ist für \betonung[b, p]{jeden} Menschen gestorben. Du musst dir selber nur die gleiche Frage stellen, die Jesus auch an Martha stellte:
        \begin{bibelbox}{SCHL}{Joh}{11:25}
            Jesus spricht zu ihr: Ich bin die Auferstehung und das Leben. Wer an mich glaubt, wird leben, auch wenn er stirbt; und jeder, der lebt und an mich glaubt, wird in Ewigkeit nicht sterben. Glaubst du das?
        \end{bibelbox}
        \betonung[b, p]{Was ist deine Antwort?}
        Das ist unser Arzt. \betonung[b, p]{\herr{} dein Arzt} Und wenn wir heute an den Tisch des Herrn treten, denken wir daran, wie der Herr unser Arzt uns von unseren Sünden geheilt hat, damit wir für die Ewigkeit vor Gott treten können.

        Zum Abschluss lesen wir noch den die Verse \bibleverse{Ps}(147:1-1)
        Amen
    \end{block}
    
    \lied[b]{Denn ich bin gewiss}.
\end{spacing}
\section{Predigt}
\textbf{Nach der Predigt}\newline
Danken für die Predigt.
\section{Abendmahl}
\begin{itemize}
    \item [] Beten für das Brot
    \item [] Beten für den Wein
\end{itemize}
\section{Abschluss}
\begin{itemize}
    \item [] \lied[b]{Bald schon kann es sein}.
    \item [] \beten[b]{Zum Abschluss}
\end{itemize}
\begin{bibelbox}{SCHL}{IIKor}{13:13}
    Die Gnade des Herrn Jesus Christus und die Liebe Gottes und die Gemeinschaft des Heiligen Geistes sei mit euch allen! Amen
\end{bibelbox}
\end{document}